\chapter{Methodology}\label{ch:methods}

\section{The original experimental record}

The \textit{Moringa} manuscript was read page by page, and every extracted item was labelled
as a measured result, an authors' interpretation, this project's commentary, or not stated;
nothing absent was reconstructed. A radiobiology review checked every extracted item against
the page images. Internal inconsistencies were recorded as printed rather than silently
corrected: the extract concentration appears as 25--50~µg/mL in the abstract, 30--50~mg/mL
in the results and above 20~µg/mL in the discussion; V79 is described variously as hamster,
murine and rodent; and the tabulated integrated areas are all positive although the plotted
optical-density changes are mostly negative for the two carcinoma lines, which cannot be
reconciled without the missing pages or the raw data.

\section{Data acquisition, harmonisation and eligibility}

Outcome, expression, identity and pathway resources were acquired under recorded
provenance, parsed into separate per-endpoint tables and harmonised on DepMap identifiers.
Download code was written only after each source had been verified; raw files are
immutable, no value was invented or silently imputed, and missingness is documented.

Cell lines were matched by stable identifiers, never by name. Cellosaurus supplied
accessions, synonyms, problem flags, and the hierarchy and origin links used to detect
relatedness. A record was eligible if it had a usable clonogenic SF2 value, a matched
baseline expression profile and a resolved identity, which gave \textbf{241 eligible
records} from 20 tissues.

\section{Endpoint and features}\label{sec:endpoint}

The single primary endpoint is \textbf{CLARIFi mSF2}, the median of the literature-reported
clonogenic surviving fractions at 2~Gy for a cell line, analysed on its raw scale.
Predictions are clipped to $[0,1]$ by a prespecified rule; outcomes are never clipped.
Crystal-violet proliferation, cell viability, ATP-based integral survival and differential
expression are kept strictly separate from clonogenic SF2 throughout the project.

Three kinds of feature were used: tissue of origin (OncotreeLineage, one-hot encoded over
fixed categories), an optional proliferation covariate, and 50 Hallmark gene-set scores
computed from baseline expression (DepMap $\log_2(\mathrm{TPM}+1)$ values of protein-coding
genes). Within each training fold, every gene $g$ is standardised with the mean $\mu_g$ and
standard deviation $\sigma_g$ of the training rows only, and the score of set $S$ for cell
line $i$ is the mean standardised expression of its member genes present in the matrix,
\begin{equation}
  s_{iS} \;=\; \frac{1}{|S^{\ast}|}\sum_{g\in S^{\ast}} \frac{x_{ig}-\mu_g}{\sigma_g},
  \qquad S^{\ast} = S \cap \{\text{measured genes}\}.
  \label{eq:score}
\end{equation}
The proliferation covariate is the combined score of the E2F-targets and G2M-checkpoint
sets.

\section{Grouping and cross-validation}\label{sec:cv}

Records of cell lines related by identity (isogenic derivatives, sublines, documented
relatedness) were merged into \textbf{237 independent biological components}; seven
components contain more than one record. Every split is grouped by component, so related
lines can never fall on both sides of a train/test boundary. Splits are stratified by
tissue: the nine tissues with at least ten components each form their own stratum and the
remaining eleven are pooled.

Validation used \textbf{ten repeats of five-fold cross-validation} with an inner five-fold
loop for hyperparameter selection. All preprocessing, scaling, gene-set scoring and tuning
took place \textbf{inside training folds only}. Each record $i$ belonging to component
$c(i)$ carries the weight
\begin{equation}
  w_i \;=\; \frac{1}{n_{c(i)}},
  \label{eq:weights}
\end{equation}
where $n_c$ is the number of records of component $c$, in every fit and every inner score,
so that a component with several records counts once. The base random seed (20261021) and
all grids are recorded in the experiment configuration. A dedicated leakage audit covered
identifier handling, component construction and split construction.

\section{Models}\label{sec:models}

Five models were prespecified, from simple to rich (Table~\ref{tab:models}), and simple
scikit-learn estimators were used before any more complex architecture. M1 is a ridge
regression on the fixed tissue categories with an unpenalised intercept, which shrinks each
tissue mean towards the overall mean. M2 is an \emph{offset model}: within each training
fold it first fits M1 and then an elastic net to the M1 residuals using the standardised
gene-set scores $\tilde{\mathbf{s}}_i$,
\begin{equation}
  \hat{y}_i^{\,\mathrm{M2}} \;=\; \hat{y}_i^{\,\mathrm{M1}} + \beta_0 +
     \tilde{\mathbf{s}}_i^{\top}\boldsymbol{\beta},
  \qquad
  \hat{\boldsymbol{\beta}} = \arg\min_{\beta_0,\boldsymbol{\beta}}
     \frac{\sum_i w_i\bigl(r_i-\beta_0-\tilde{\mathbf{s}}_i^{\top}\boldsymbol{\beta}\bigr)^2}
          {2\sum_i w_i}
     + \alpha\rho\,\lVert\boldsymbol{\beta}\rVert_1
     + \frac{\alpha(1-\rho)}{2}\lVert\boldsymbol{\beta}\rVert_2^2 ,
  \label{eq:offset}
\end{equation}
with $r_i = y_i-\hat{y}_i^{\,\mathrm{M1}}$. Because M2 contains M1 as an offset, the
comparison of the two asks specifically what expression adds beyond tissue. M1P and M2P are
the corresponding models with the proliferation covariate added to the baseline; M2P uses
the remaining 48 Hallmark scores in its residual model. The inner loop searched
$\alpha\in\{0.1, 0.01, 0.001\}$, $\rho\in\{0.9, 0.5, 0.1\}$ and the baseline ridge penalty
$\in\{100, 10, 1, 0.1\}$, 36 combinations for M2.

\begin{table}[tbp]
  \caption{The five prespecified models of the primary benchmark.}\label{tab:models}
  \small
  \begin{tabular}{@{}l L{0.50\linewidth} L{0.30\linewidth}@{}}
    \toprule
    \thead{Model} & \thead{Features} & \thead{Role} \\
    \midrule
    M0 & None (training-set mean) & Floor \\
    M1 & Tissue of origin & Primary comparator \\
    M1P & Tissue and proliferation score & Sensitivity comparator \\
    \textbf{M2} & \textbf{Tissue and 50 Hallmark expression scores} & \textbf{Primary model} \\
    M2P & Tissue, proliferation and the other 48 Hallmark scores & Sensitivity model \\
    \bottomrule
  \end{tabular}
\end{table}

Two further models were configured for the exploratory layer
(Chapter~\ref{ch:robustness}): M2X, an elastic net on the 50 scores with the global mean as
offset and no tissue information, used only for the leave-one-tissue-out analysis; and M3,
a random forest on one-hot tissue and the 50 scores. M3 was added after the primary results
had been seen and therefore cannot replace M2 as the primary benchmark, whatever its score.

\section{Performance measures and the primary estimand}\label{sec:metrics}

For a repeat $r$ with out-of-fold predictions $\hat{y}_{ir}$, the component-weighted error
measures are
\begin{align}
  \mathrm{MAE}_r &= \frac{\sum_i w_i\,\lvert y_i-\hat{y}_{ir}\rvert}{\sum_i w_i}, &
  \mathrm{RMSE}_r &= \Biggl(\frac{\sum_i w_i\,(y_i-\hat{y}_{ir})^2}{\sum_i w_i}\Biggr)^{1/2},
  \label{eq:mae}\\
  R^2_{\mathrm{oos},r} &= 1-\frac{\sum_i w_i\,(y_i-\hat{y}_{ir})^2}
                                  {\sum_i w_i\,(y_i-\bar{y}_w)^2}, &
  \bar{y}_w &= \frac{\sum_i w_i y_i}{\sum_i w_i},
  \label{eq:r2}
\end{align}
and each reported value is the mean over the $R=10$ repeats. $R^2_{\mathrm{oos}}=0$
therefore corresponds to predicting the cohort mean, and negative values are worse than
that. Calibration slope and intercept come from a weighted regression of observed on
predicted values.

The primary estimand is the paired, component-weighted difference in mean absolute error
between M1 and M2. With the record-level difference
$d_{ir}=\lvert y_i-\hat{y}^{\,\mathrm{M1}}_{ir}\rvert-\lvert y_i-\hat{y}^{\,\mathrm{M2}}_{ir}\rvert$,
the component means $D_{cr}=n_c^{-1}\sum_{i\in c} d_{ir}$ and the $C=237$ components,
\begin{equation}
  \Delta\mathrm{MAE} \;=\; \frac{1}{R}\sum_{r=1}^{R}\,\frac{1}{C}\sum_{c=1}^{C} D_{cr},
  \label{eq:dmae}
\end{equation}
so a positive value favours the richer model. MAE, RMSE and $R^2_{\mathrm{oos}}$ are
reported for every model, and no metric was chosen after the results were seen.

\section{Explainability}\label{sec:explain-methods}

For each of the 50 real-data fits of M2 (ten repeats by five folds), grouped permutation
importance was computed on the \emph{outer-test rows only}, using the saved model, with 20
permutations per feature; nothing was refitted. For gene set $S$ in fit $f$,
\begin{equation}
  I_{S,f} \;=\; \frac{1}{20}\sum_{p=1}^{20}
     \Bigl[\mathrm{MAE}_f\bigl(\text{score } S \text{ permuted across components, draw } p\bigr)
           -\mathrm{MAE}_f\Bigr],
  \label{eq:importance}
\end{equation}
the increase in component-weighted held-out error, in mSF2 units, when the information in
that score is destroyed. Coefficient sign consistency across fits and the Spearman
correlation between gene-set scores were also computed, so that shared signal between
correlated sets is visible rather than hidden.

\section{Pathway prioritisation}\label{sec:pathway-methods}

Gene sets were assigned to tiers by rules frozen before the run. \emph{Tier~C} requires that
importance exceeded zero in at least half the fits and that the share of fits with a zero
coefficient was at most 0.50. \emph{Tier~B} additionally requires sign consistency and a
positive-importance share of at least 0.75 and significance against the permuted-outcome
null; \emph{tier~A} requires both at least 0.90 and significance against both null models.
\textbf{Both higher tiers also require a formally demonstrated prediction gain and at least
1,000 completed draws of each null.} Neither exists in this run, so tiers A and B were
unreachable by construction, and the rules were not altered to produce a more impressive
ranking. Scores correlated at $\lvert\rho\rvert\ge 0.80$ are flagged as sharing signal.

\section{Cell-line triage}\label{sec:triage-methods}

Each measured line was characterised by the absolute error of its repeat-averaged
out-of-fold prediction under M1 and M2, the spread of its prediction across repeats, and
the disagreement between the two models, each expressed as a cohort percentile (the share
of lines with a strictly smaller value, plus one, divided by the number of lines). A line
reaches \emph{tier~A} when its error is at or above the 95th percentile under both models,
\emph{tier~B} at or above the 90th percentile under both, and \emph{tier~C} at or above the
90th percentile under either. Technical flags --- stranded RNA profile, related lines in the
cohort, a Cellosaurus problem flag on a relative, a small tissue group, unstable prediction
and strong model disagreement --- are recorded separately from the tier, so that a large
residual caused by a data or identity problem is never presented as a biological
discrepancy.

\section{Curation of \textit{Moringa} constituent evidence}\label{sec:curation-methods}

Constituents and their biology were curated by literature search and reading only, into two
validated evidence files. Only \emph{primary} studies count, so a review citing an earlier
study is never independent confirmation. Only studies of the \textbf{purified compound} can
score; a study that dosed an extract is recorded with its score contribution disabled. Only
\emph{ionising} radiation qualifies as radiation evidence, which excludes the large UV
photoprotection literature. Nanoparticle and metal-conjugate formulations are excluded
because the carrier confounds the compound, and chemically distinct derivatives are never
credited to the parent compound. The \textbf{extraction solvent of every report is recorded
verbatim}, and each compound is classified as documented in a dichloromethane extract,
documented in seed or leaf with another solvent, or indirect and unconfirmed. Every DOI was
verified against Europe PMC or Crossref and every identifier against PubChem, with the
verification route and reading depth stored per record.

\section{Candidate scoring and tiers}\label{sec:scoring}

Each compound receives four integer scores, fixed before any score was computed:
\begin{equation}
  T \;=\; \underbrace{D_1}_{0\text{--}2}+\underbrace{D_2}_{0\text{--}3}
         +\underbrace{D_3}_{0\text{--}2}+\underbrace{D_4}_{0\text{--}1}\;\in\{0,\dots,8\},
  \label{eq:priority}
\end{equation}
where $D_1$ is the strength of verified \textit{Moringa} constituent evidence (2 if at least
one primary seed or leaf report has structural confirmation by isolation and NMR or by a
reference standard; 1 if identification rests on a spectral-library match or is not stated;
0 if no primary seed or leaf report exists); $D_2$ the relevance to radiation-response
biology (3 for ionising-radiation evidence from at least two independent research groups,
2 from one group, 1 if only a relevant process was modulated without radiation); $D_3$
independent source support (2 for at least three distinct research groups across both
evidence files, 1 for two); and $D_4$ feasibility of follow-up (1 if the identifier is
verified and at least one scored study reports a usable concentration or dose). Research
groups are counted once across both evidence files, so repeated work from one laboratory is
one source, and \textbf{chemical-class analogy earns no points}.

The tiers are A ($T\ge 6$, $D_1\ge 1$, $D_2\ge 2$), B ($T\ge 4$, $D_1\ge 1$, $D_2\ge 1$) and C
($T\ge 2$, $D_1\ge 1$); compounds below these thresholds are not prioritised. Two safeguards
override the score. A compound with no primary \textit{Moringa} identification ($D_1=0$) is
\emph{not assessable} and never enters the shortlist. Evidence that a compound protects
\emph{tumour} cells from radiation, or is itself genotoxic, is a liability that is displayed
and caps the tier at C. Extraction-specific evidence is reported beside the score and never
folded into it.

Model-associated pathways, literature-supported processes and their overlap are kept apart.
An overlap between the genes a published experiment measured and the members of a
model-ranked gene set marks a hypothesis to be tested; it is not evidence that the compound
modifies the gene set, the pathway or clonogenic SF2.

\section{Reproducibility and quality assurance}\label{sec:repro}

The codebase uses Python~3.11 or later with a \path{src/} layout, type hints, small pure
functions and YAML configuration with recorded seeds and no machine-specific paths. Every
analytical output carries a run record with the git commit and a flag for uncommitted
changes, configuration and file checksums, package versions and UTC timestamps, and
analysis runs refuse to start from an uncommitted working tree. Result folders and figures
are write-once, so a published number cannot be silently overwritten. Specialist reviews
covering radiobiology, radiogenomic data curation, machine-learning methodology, statistics,
software architecture, reproducibility, figure design and literature were recorded at each
decision point in the project's decision log, together with any unresolved disagreement.
